\documentclass[10pt,twocolumn]{article}

% ── Packages ──────────────────────────────────────────────────────────────
\usepackage[margin=0.75in,top=0.9in,bottom=0.9in]{geometry}
\usepackage{times}
\usepackage{microtype}
\usepackage{amsmath,amssymb}
\usepackage{booktabs}
\usepackage{graphicx}
\usepackage{xcolor}
\usepackage{hyperref}
\usepackage{enumitem}
\usepackage{multicol}
\usepackage{caption}
\usepackage{authblk}
\usepackage{abstract}
\usepackage{titlesec}
\usepackage{parskip}
\usepackage{array}

% ── Hyperref setup ─────────────────────────────────────────────────────────
\hypersetup{
  colorlinks=true,
  linkcolor=blue!60!black,
  citecolor=blue!60!black,
  urlcolor=blue!60!black,
}

% ── Section title formatting ───────────────────────────────────────────────
\titleformat{\section}{\normalfont\large\bfseries}{\thesection.}{0.5em}{}
\titleformat{\subsection}{\normalfont\normalsize\bfseries}{\thesubsection.}{0.5em}{}
\titlespacing*{\section}{0pt}{8pt}{4pt}
\titlespacing*{\subsection}{0pt}{6pt}{3pt}

% ── Abstract box ───────────────────────────────────────────────────────────
\renewcommand{\abstractnamefont}{\normalfont\bfseries}
\setlength{\absleftindent}{0pt}
\setlength{\absrightindent}{0pt}

% ── Compact lists ──────────────────────────────────────────────────────────
\setlist[itemize]{topsep=2pt,itemsep=1pt,parsep=0pt,leftmargin=*}
\setlist[enumerate]{topsep=2pt,itemsep=1pt,parsep=0pt,leftmargin=*}

% ── Table column helpers ───────────────────────────────────────────────────
\newcolumntype{L}[1]{>{\raggedright\arraybackslash}p{#1}}
\newcolumntype{C}[1]{>{\centering\arraybackslash}p{#1}}
\newcolumntype{R}[1]{>{\raggedleft\arraybackslash}p{#1}}

% ─────────────────────────────────────────────────────────────────────────
\begin{document}

% ── Title & Authors ───────────────────────────────────────────────────────
\title{\textbf{Grounded Vision-Language Agents\\for Instruction Following}}

\author[1]{Karthik Reddy Jannupalli}
\author[2]{Sri Sai Prasanna Konyala}
\author[3]{Sai Syama Srikar Parimi}
\author[4]{Prajeeth Bharadwaj Kittoor Muralidhar}
\author[5]{Satya Bhargav Maddipati}
\author[6]{Sanjana Reji Kallingal}

\affil[ ]{\small
  \texttt{jannupalli@wisc.edu}\quad
  \texttt{skonyala@wisc.edu}\quad
  \texttt{sparimi2@wisc.edu}\quad
  \texttt{kittoormural@wisc.edu}\quad
  \texttt{sbmaddipati@wisc.edu}\quad
  \texttt{kallingal@wisc.edu}\\[2pt]
  Department of Computer Sciences, University of Wisconsin--Madison\\
  CS\,639: Introduction to Foundation Models, Spring 2026
}

\date{}
\maketitle
\thispagestyle{empty}

% ── Abstract ──────────────────────────────────────────────────────────────
\begin{abstract}
\noindent
We study grounded instruction following in visual environments: agents that must not only
understand a natural-language command but locate and act on the correct target in a screenshot.
We build a complete evaluation pipeline with three agent variants and benchmark them on
ScienceQA and a synthetic grounded-UI corpus. Our results confirm three hypotheses:
(H1)~visual grounding is necessary, adding 23 percentage points over a text-only ReAct
baseline; (H2)~an iterative Observe--Reason--Act (ORA) loop with per-step visual re-encoding
provides a further 10\,pp gain over single-shot VLM prompting; and (H3)~LoRA fine-tuning on
200 synthetic examples yields 92.0\% accuracy on the training domain with zero catastrophic
forgetting on ScienceQA (57.5\% vs.\ 56.5\% baseline). We additionally report a practical
finding: label masking---excluding prompt and image patch tokens from the cross-entropy
loss---is critical for multimodal LoRA training and reduces training loss by a factor of 888$\times$.
\end{abstract}

\vspace{-4pt}
\noindent\rule{\linewidth}{0.4pt}
\vspace{2pt}

% ── 1. Introduction ───────────────────────────────────────────────────────
\section{Introduction}

Autonomous agents that can follow natural-language instructions in visual environments have broad
potential: from assistive tools for users with motor impairments, to automated testing of
web interfaces, to general-purpose digital assistants. The core challenge is
\emph{grounded instruction following}---the agent must not only understand what to do, but
locate the correct target in a visual scene (a button, a diagram element, a form field) and
produce a precise spatial action.

Recent large vision-language models (VLMs) such as LLaVA~\cite{liu2024llava} have demonstrated
strong zero-shot performance on visual QA, but their deployment as \emph{action-taking agents}
in multi-step grounded tasks is less studied. Two open questions motivate this work:

\begin{enumerate}
  \item Does adding vision to a reasoning agent meaningfully improve performance on spatially
        grounded tasks, or is a capable text-only model sufficient?
  \item Can a structured Observe$\to$Reason$\to$Act (ORA) loop with per-step visual re-encoding
        outperform single-shot VLM prompting---and can lightweight LoRA fine-tuning on a small
        synthetic corpus push performance further without degrading general capability?
\end{enumerate}

We build a complete pipeline, three agent variants, and run eight controlled experiments
comparing them on two benchmarks. All three hypotheses are confirmed.

% ── 2. Research Hypotheses ────────────────────────────────────────────────
\section{Research Hypotheses}

We evaluate three nested hypotheses, each isolating a single design choice:

\textbf{H1 -- Visual Grounding.}
A multimodal agent (LLaVA-1.6-7B) will substantially outperform a text-only agent
(Mistral-7B + ReAct) on spatially grounded tasks, because correct action prediction requires
interpreting the visual layout, not just the textual task description.

\textbf{H2 -- Structured ORA Reasoning.}
An iterative Observe-Reason-Act loop that re-encodes the screenshot at every step will
outperform a single-shot VLM call, because the agent can condition each action on the
\emph{current} visual state rather than a stale textual summary.

\textbf{H3 -- LoRA Adaptation (stretch goal).}
Fine-tuning ORA-LLaVA with a LoRA adapter~\cite{hu2022lora} on a small (200-example) grounded
synthetic corpus will improve task completion on the training domain without causing catastrophic
forgetting on general visual QA, because LoRA updates only 0.21\% of model parameters.

% ── 3. Related Work ───────────────────────────────────────────────────────
\section{Related Work}

\textbf{Vision-Language Models.}
LLaVA-1.6~\cite{liu2024llava} extends instruction-tuned LLMs with a visual encoder connected
via a projection layer, enabling visual QA and instruction following. We use LLaVA-1.6-Mistral-7B
as the backbone for H2 and H3. Unlike standard VQA benchmarks, our evaluation focuses on
\emph{action prediction} with precise spatial coordinates---a harder test of visual grounding.

\textbf{ReAct.}
Yao et al.~\cite{yao2023react} introduced ReAct, which interleaves natural-language reasoning
traces and environment actions. We use Mistral-7B-Instruct~\cite{jiang2023mistral} with ReAct
as our text-only baseline (H1); our ORA loop extends ReAct to multimodal inputs.

\textbf{Mind2Web.}
Deng et al.~\cite{deng2023mind2web} released Mind2Web, a web navigation benchmark with real
browser interactions. Our synthetic task design borrows Mind2Web's action taxonomy (click, type,
scroll), adapted to a controlled programmatic visual setting.

\textbf{ScienceQA.}
Lu et al.~\cite{lu2022scienceqa} introduced ScienceQA, a multi-modal multiple-choice benchmark
covering science topics with diagrams. We use it as a held-out general-capability benchmark to
measure forgetting after LoRA fine-tuning.

\textbf{Parameter-Efficient Fine-Tuning.}
Hu et al.~\cite{hu2022lora} proposed LoRA, which inserts low-rank update matrices into frozen
transformer layers, drastically reducing trainable parameters. We use the PEFT
library~\cite{mangrulkar2022peft}. A key finding is that \emph{label masking}---excluding prompt
and image tokens from the cross-entropy loss---is critical when applying LoRA to multimodal
inputs (Section~\ref{sec:lora}).

\textbf{VLM-Based Web Agents.}
Zheng et al.~\cite{zheng2024seeact} demonstrated SeeAct using GPT-4V for web agent tasks.
Our work differs in using an open 7B model with LoRA fine-tuning, making it accessible
for researchers without proprietary API access.

% ── 4. Methodology ───────────────────────────────────────────────────────
\section{Methodology}

\subsection{Agent Designs}

We implement three agents sharing the same evaluation harness.

\textbf{ReAct-Mistral (Baseline 1).}
A text-only agent using Mistral-7B-Instruct with the standard ReAct prompt. The agent receives
the task instruction and a textual description of the visual state; it produces a
Thought\,+\,Action at each step. This isolates the effect of removing vision while keeping the
reasoning loop.

\textbf{Single-Shot LLaVA (Baseline 2).}
A single-call VLM agent using LLaVA-1.6-Mistral-7B (4-bit quantized). The agent receives the
task instruction and screenshot in a single forward pass and produces one action. This isolates
the effect of removing the iterative loop while keeping vision.

\textbf{ORA-LLaVA (Our Method, H2).}
The same LLaVA backbone in an iterative loop. At each step, the current screenshot is
re-encoded through the full vision tower, grounding the action prediction in the \emph{current}
visual state. This is the key distinction from ReAct: the visual observation is never
summarized into text---raw image features enter the transformer at every step.

\subsection{LoRA Fine-Tuning (H3)}
\label{sec:lora}

We fine-tune ORA-LLaVA with a LoRA adapter ($r{=}16$, $\alpha{=}32$, dropout $= 0.05$,
$\text{lr} = 2{\times}10^{-4}$, 3 epochs) on 200 synthetic grounded tasks. Each training
example is a (screenshot, instruction, action) triple. The assistant response (a JSON action
string) is the training target.

\textbf{Critical implementation finding.}
Our initial training run plateaued at loss ${\approx}4.0$ despite correct hyperparameters.
Diagnosis: the loss was computed over \emph{all} input tokens, including approximately 1,500
image patch tokens whose values are unpredictable. The fix is to mask all prompt tokens
(labels$[{:}\text{prompt\_len}] = -100$), so only the ${\approx}30$ assistant response tokens
contribute to the cross-entropy loss. After this fix, loss converged to 0.0045 at step 70---an
$888{\times}$ improvement---and evaluation accuracy on the training domain jumped from 15.5\%
to 92.0\%.

Trainable parameters: 15,990,784 / 7,582,738,432 (0.21\%). Training time: ${\approx}17$~min
on a single NVIDIA A100 40\,GB.

\subsection{Datasets}

\textbf{ScienceQA ($n{=}200$).}
We sample 200 test-split examples from ScienceQA involving diagrams. Used as the primary
benchmark for H1/H2 and as a held-out general-capability test for H3.

\textbf{Synthetic Grounded Tasks ($n{=}200$).}
We programmatically generate 200 PNG screenshots of simulated UI elements (login forms,
navigation bars, charts, calendars, etc.) using Pillow, with 30 distinct image generators
covering diverse visual layouts. Each task pairs an instruction (e.g.,
``\textit{click the Submit button}'') with a ground-truth click action and pixel coordinates.
A separate 200-example split is used for LoRA training.

\subsection{Evaluation Protocol}

A task is counted as a \emph{success} if the agent produces a correct action within the step
budget. We record (i) task completion rate (primary metric), (ii) mean steps per task
(efficiency), and (iii) error breakdown across four categories:
\emph{visual misgrounding}, \emph{reasoning error}, \emph{action parse failure},
and \emph{truncated} (step budget exceeded without success).

% ── 5. Results ────────────────────────────────────────────────────────────
\section{Results}

\subsection{Main Results}

Table~\ref{tab:main} summarizes task completion rates for all eight runs.

\begin{table}[ht]
\centering
\caption{Task completion rate (\%) across all agents and datasets.
  $\dagger$~Baseline synthetic evaluation used $n{=}46$.}
\label{tab:main}
\small
\begin{tabular}{@{}lcc@{}}
\toprule
\textbf{Agent} & \textbf{ScienceQA} & \textbf{Synthetic}\\
               & $n=200$            & $n=200$\\
\midrule
ReAct-Mistral          & 23.5\%  & 13.0\%$^\dagger$\\
Single-Shot LLaVA      & 46.5\%  & 37.0\%$^\dagger$\\
ORA-LLaVA              & 56.5\%  & 41.3\%$^\dagger$\\
\midrule
\textbf{ORA-LLaVA + LoRA} & \textbf{57.5\%} & \textbf{92.0\%}\\
\bottomrule
\end{tabular}
\end{table}

\begin{table}[ht]
\centering
\caption{Detailed error breakdown (\% of $n$ tasks).
  VM = visual misgrounding; RE = reasoning error; PF = parse failure; TR = truncated.}
\label{tab:errors}
\small
\setlength{\tabcolsep}{4pt}
\begin{tabular}{@{}llccccc@{}}
\toprule
\textbf{Agent} & \textbf{Dataset} & \textbf{Succ.} & \textbf{VM} & \textbf{RE} & \textbf{PF} & \textbf{TR}\\
\midrule
ReAct-Mistral  & ScienceQA & 23.5 & 52.5 & 11.0 & 12.5 & 0.5\\
ReAct-Mistral  & Synthetic  & 13.0 & 15.2 & 41.3 & 17.4 & 13.0\\
Single-Shot    & ScienceQA & 46.5 & 26.5 & 14.0 & 13.0 & 0.0\\
Single-Shot    & Synthetic  & 37.0 & 41.3 & 10.9 & 10.9 & 0.0\\
ORA-LLaVA      & ScienceQA & 56.5 & 33.0 &  7.0 &  0.0 & 3.5\\
ORA-LLaVA      & Synthetic  & 41.3 & 17.4 &  0.0 &  4.3 & 37.0\\
\midrule
\textbf{+LoRA} & ScienceQA & \textbf{57.5} & 42.5 & 0.0 & 0.0 & 0.0\\
\textbf{+LoRA} & Synthetic  & \textbf{92.0} &  1.5 & 2.0 & 0.0 & 4.5\\
\bottomrule
\end{tabular}
\end{table}

\subsection{H1: Visual Grounding Confirmed}

Single-Shot LLaVA (46.5\%) outperforms ReAct-Mistral (23.5\%) by +23\,pp on ScienceQA and
+24\,pp on Synthetic (37.0\% vs.\ 13.0\%). The text-only model's primary failure is visual
misgrounding (52.5\% of failures on ScienceQA), confirming that spatial layout information
is not recoverable from text alone.

\subsection{H2: ORA Loop Confirmed}

ORA-LLaVA (56.5\%) outperforms Single-Shot (46.5\%) by +10\,pp on ScienceQA with the same
LLaVA backbone. The iterative re-encoding reduces reasoning errors by 50\% (28$\to$14) and
eliminates parse failures (26$\to$0), confirming that per-step visual grounding keeps the
agent on track.

\subsection{H3: LoRA Adaptation Strongly Confirmed}

After fine-tuning with the label-masking fix:
\begin{itemize}
  \item \textbf{Synthetic: 92.0\%} (+50.7\,pp over ORA baseline; +122\% relative).
        184/200 tasks solved in a single step; truncations fall from 17 to 9.
  \item \textbf{ScienceQA: 57.5\%}---essentially unchanged from the 56.5\% baseline,
        confirming zero catastrophic forgetting despite training exclusively on UI tasks.
  \item Reasoning errors and truncations on ScienceQA drop to \emph{zero}, indicating
        the adapter improved the model's output discipline beyond the training domain.
\end{itemize}
Mean steps on ScienceQA fell from 1.23 (ORA) to 1.00 (LoRA), indicating greater decisiveness.

\subsection{Error Analysis}

Visual misgrounding is the dominant failure mode across all models (Table~\ref{tab:errors}).
After LoRA fine-tuning, synthetic visual misgrounding drops to 1.5\% (3/200), but ScienceQA
visual misgrounding \emph{increases} to 42.5\% (85/200)---rising as truncation and reasoning
errors are eliminated. This indicates the model's residual errors are purely spatial and cannot
be resolved by response-token fine-tuning alone; future work should unfreeze the vision tower.

% ── 6. Conclusions ───────────────────────────────────────────────────────
\section{Conclusions and Future Work}

This project provides three empirical findings. First, \textbf{vision is necessary}: text-only
ReAct achieves only 23.5\% on visually grounded tasks, while the same loop with a VLM backbone
reaches 56.5\%. Second, \textbf{iterative visual re-encoding helps}: the ORA loop provides a
consistent +10\,pp over single-shot prompting, validating the design choice of re-encoding the
screenshot at every step. Third, \textbf{LoRA fine-tuning on 200 synthetic examples is highly
effective}: 92\% accuracy on the training domain with zero forgetting on the held-out benchmark,
training only 0.21\% of model parameters in 17 minutes.

An unexpected practical contribution is the identification of \textbf{label masking as a
prerequisite for multimodal LoRA training}. Without masking prompt and image patch tokens from
the cross-entropy loss, training converges to a degenerate solution (loss plateau at ${\sim}4.0$,
eval accuracy 4.5\%). This finding is not consistently documented in fine-tuning tutorials and
represents a critical implementation requirement for any multimodal PEFT work.

\textbf{Future directions:}
(1)~Scale the synthetic dataset to 2{,}000 tasks with harder layouts (occlusion, dense UIs,
small targets) to test whether LoRA generalizes beyond memorization.
(2)~Unfreeze the vision tower in a second LoRA pass to address persistent visual misgrounding.
(3)~Evaluate on real web navigation benchmarks (Mind2Web, WebArena) for out-of-distribution
generalization.
(4)~Explore chain-of-thought visual description as an additional observation channel before
action prediction.

% ── References ───────────────────────────────────────────────────────────
\section*{References}
{\small
\begin{enumerate}[label={[\arabic*]},leftmargin=*,topsep=4pt,itemsep=2pt]

  \item\label{ref:deng2023} X.~Deng, Y.~Gu, B.~Zheng, S.~Chen, S.~Stevens, B.~Wang,
    H.~Sun, and Y.~Su,
    ``Mind2Web: Towards a generalist agent for the web,''
    in \emph{Proc.\ NeurIPS}, 2023.

  \item\label{ref:hu2022} E.~J.~Hu, Y.~Shen, P.~Wallis, Z.~Allen-Zhu, Y.~Li, S.~Wang,
    L.~Wang, and W.~Chen,
    ``LoRA: Low-rank adaptation of large language models,''
    in \emph{Proc.\ ICLR}, 2022.

  \item\label{ref:jiang2023} A.~Q.~Jiang, A.~Sablayrolles, A.~Mensch, C.~Bamford,
    D.~S.~Chaplot, D.~de~las~Casas, F.~Bressand, G.~Lengyel, G.~Lample,
    L.~Saulnier, et al.,
    ``Mistral 7B,''
    \emph{arXiv:2310.06825}, 2023.

  \item\label{ref:liu2024} H.~Liu, C.~Li, Y.~Li, and Y.~J.~Lee,
    ``Improved baselines with visual instruction tuning (LLaVA-1.5),''
    in \emph{Proc.\ CVPR}, 2024.

  \item\label{ref:lu2022} P.~Lu, S.~Mishra, T.~Xia, L.~Qiu, K.-W.~Chang, S.-C.~Zhu,
    O.~Tafjord, P.~Clark, and A.~Kalyan,
    ``Learn to explain: Multimodal reasoning via thought chains for science question answering,''
    in \emph{Proc.\ NeurIPS}, 2022.

  \item\label{ref:mangrulkar2022} S.~Mangrulkar, S.~Gugger, L.~Debut, Y.~Belkada,
    S.~Paul, and B.~Bossan,
    ``PEFT: State-of-the-art parameter-efficient fine-tuning methods,''
    \url{https://github.com/huggingface/peft}, 2022.

  \item\label{ref:yao2023} S.~Yao, J.~Zhao, D.~Yu, N.~Du, I.~Shafran,
    K.~R.~Narasimhan, and Y.~Cao,
    ``ReAct: Synergizing reasoning and acting in language models,''
    in \emph{Proc.\ ICLR}, 2023.

  \item\label{ref:zheng2024} B.~Zheng, B.~Gou, J.~Kil, H.~Sun, and Y.~Su,
    ``GPT-4V(ision) is a generalist web agent, if grounded,''
    \emph{arXiv:2401.01614}, 2024.

\end{enumerate}
}

% ── BibTeX keys (for \cite commands above) ────────────────────────────────
% liu2024llava   -> [4]
% hu2022lora     -> [2]
% yao2023react   -> [7]
% jiang2023mistral -> [3]
% deng2023mind2web -> [1]
% lu2022scienceqa  -> [5]
% mangrulkar2022peft -> [6]
% zheng2024seeact  -> [8]

\end{document}
